\chapter{Référentiels et Critères d'Évaluation}
\label{chap:refaudit}

Un auditeur a besoin d'un "mètre étalon" pour mesurer la conformité. Ce mètre est le \textbf{référentiel}. Sans référentiel, l'audit n'est qu'une simple opinion.

\section{La Structure Commune (High Level Structure)}
% ------------------------------------------------------------------------------

Depuis 2012, les normes de management (ISO 9001, ISO 27001, ISO 22301) partagent une structure commune pour faciliter leur intégration.

\begin{boxdef}
\textbf{High Level Structure (HLS) :} Structure identique en 10 chapitres pour toutes les normes ISO de management. Cela permet d'aligner Qualité, Sécurité et Environnement dans un Système de Management Intégré (SMI).
\end{boxdef}

Les chapitres 4 à 10 sont les chapitres "exigentiels" (ceux sur lesquels on audite) :
\begin{itemize}
    \item \textbf{Chapitre 4 :} Contexte de l'organisme.
    \item \textbf{Chapitre 5 :} Leadership (Engagement de la direction).
    \item \textbf{Chapitre 6 :} Planification (Analyse de risque).
    \item \textbf{Chapitre 7 :} Support (Ressources, Compétences, Sensibilisation).
    \item \textbf{Chapitre 8 :} Réalisation (Opérations, Plan d'action).
    \item \textbf{Chapitre 9 :} Évaluation des performances (Audit interne, Indicateurs).
    \item \textbf{Chapitre 10 :} Amélioration (Non-conformités, Actions correctives).
\end{itemize}

% ------------------------------------------------------------------------------
\section{ISO 27001 : La Référence Sécurité (SMSI)}
% ------------------------------------------------------------------------------

\subsection{Présentation}

\begin{boxdef}
\textbf{ISO 27001} : Norme internationale définissant les exigences pour établir, mettre en œuvre, maintenir et améliorer un Système de Management de la Sécurité de l'Information (SMSI).
\end{boxdef}

C'est la norme "système" : elle dit \textit{comment} manager la sécurité (PDCA).

\subsection{L'Annexe A : Mesures de Sécurité}

C'est la partie la plus connue des techniciens. L'Annexe A contient une liste de \textbf{mesures (contrôles)} à implémenter pour réduire les risques.

\begin{boxlizbiz}
\textbf{Structure de l'Annexe A (ISO 27001:2022) chez LiZbiZ :}
\begin{itemize}
    \item \textbf{A.5 :} Politiques de sécurité.
    \item \textbf{A.6 :} Organisation de la sécurité (Rôles RACI).
    \item \textbf{A.8 :} Sécurité liée aux ressources humaines (Embauche, Départ).
    \item \textbf{A.9 :} Contrôle d'accès (Mots de passe, Privilèges).
    \item \textbf{A.12 :} Sécurité des opérations (Logs, Sauvegardes, Malware).
\end{itemize}
\end{boxlizbiz}

\subsection{Notion d'Exigence vs Recommandation}

En audit, la nuance est importante :
\begin{itemize}
    \item \textbf{Exigence ("Shall") :} Obligatoire pour la certification.
    \item \textbf{Recommandation ("Should") :} Bonne pratique, mais non bloquante.
\end{itemize}

\begin{boxalert}
Dans l'ISO 27001, le corps de la norme contient des \textbf{exigences}. L'Annexe A contient des \textbf{mesures} qui sont obligatoires \textit{si} elles ont été sélectionnées dans la Déclaration d'Applicabilité (SOA - Statement of Applicability).
\end{boxalert}

% ------------------------------------------------------------------------------
\section{ISO 9001 : La Référence Qualité}
% ------------------------------------------------------------------------------

\subsection{Présentation}

\begin{boxdef}
\textbf{ISO 9001} : Norme internationale relative aux systèmes de management de la qualité. Elle vise la satisfaction du client et l'amélioration des processus.
\end{boxdef}

Pourquoi l'inclure dans un cours Cybersécurité ? Parce que la sécurité est un composant de la qualité de service.

\subsection{Approche Processus}

L'ISO 9001 exige que l'entreprise cartographie ses processus.
\begin{itemize}
    \item \textbf{Processus de Direction :} Pilotage.
    \item \textbf{Processus de Réalisation :} Création de valeur (Vente, Livraison).
    \item \textbf{Processus Support :} Support (RH, IT, Comptabilité).
\end{itemize}

\begin{boxlizbiz}
\textbf{Lien LiZbiZ :}
Un client commande un PC (Processus Réalisation). Si le SI est piraté et les données perdues, la qualité du service est dégradée. Auditer la qualité du processus "Livraison" implique de vérifier la disponibilité du SI.
\end{boxlizbiz}

% ------------------------------------------------------------------------------
\section{COBIT : La Référence Gouvernance IT}
% ------------------------------------------------------------------------------

\subsection{Présentation}

\begin{boxdef}
\textbf{COBIT} (Control Objectives for Information and Related Technologies) : Framework de bonnes pratiques pour la gouvernance et le management des SI, développé par l'ISACA.
\end{boxdef}

Contrairement à l'ISO 27001 (focus Sécurité), COBIT couvre tout le SI : stratégie, valeur, qualité, risque, sécurité.

\subsection{Les Domaines de COBIT}

COBIT divise le management de l'IT en 4 domaines principaux :
\begin{enumerate}
    \item \textbf{Planifier et Organiser (PO) :} Stratégie, Architecture.
    \item \textbf{Acquérir et Implémenter (AI) :} Projets, Développement.
    \item \textbf{Délivrer et Servir (DS) :} Exploitation, Support utilisateur, Sécurité.
    \item \textbf{Surveiller et Évaluer (ME) :} Audit, Performance, Conformité.
\end{enumerate}

\begin{boxlizbiz}
\textbf{Utilisation pour le Management (Filière Management SI) :}
L'auditeur utilise COBIT pour auditer la maturité de la DSI de LiZbiZ.
Exemple : Le processus "DS5 - Assurer la sécurité des services" de COBIT complète l'ISO 27001 en ajoutant une vision management (budget, rôles, indicateurs).
\end{boxlizbiz}

% ------------------------------------------------------------------------------
\section{Cas LiZbiZ : Sélection des Critères d'Audit}
% ------------------------------------------------------------------------------

\subsection{Déclaration d'Applicabilité (SOA)}

LiZbiZ ne peut pas (et ne doit pas) appliquer les 93 mesures de l'Annexe A de l'ISO 27001. L'entreprise doit rédiger une \textbf{Déclaration d'Applicabilité}.

\begin{boxdef}
\textbf{SOA (Statement of Applicability) :} Document listant les mesures de l'Annexe A ISO 27001 applicables à l'entreprise, avec une justification pour celles qui sont exclues.
\end{boxdef}

\subsection{Travaux Pratiques : Construction de la Grille d'Audit}

\begin{boxlizbiz}
\textbf{Mission :} Préparer la grille d'audit pour la section "Contrôle d'Accès".

En vous basant sur les risques identifiés dans le Module 1 (Risque Ransomware, Risque Interne), sélectionnez 4 mesures pertinentes de l'Annexe A de l'ISO 27001 pour LiZbiZ.

\textbf{Exemple de grille à construire :}
\begin{table}[htbp]
\centering
\small
\begin{tabular}{p{2cm}p{6cm}p{4cm}}
    \toprule
    \textbf{Référence} & \textbf{Exigence} & \textbf{Moyen de vérification (Audit)} \\
    \midrule
    A.5.15 & Contrôle des accès logiques. & Extraction AD, Test de robustesse mdp. \\
    A.5.16 & Gestion des identités. & Procédure écriture, Test départ salarié. \\
    A.8.2 & Privilèges d'accès. & Liste des "Admin", Vérification need-to-know. \\
    A.8.3 & Retrait des droits d'accès. & Comparatif liste salariés actifs vs comptes IT. \\
    \bottomrule
\end{tabular}
\end{table}
\end{boxlizbiz}

% ==============================================================================
% Fichier : 03_module_audit.tex (Extrait Chapitre 4)
% Description : Application Pratique - Réalisation de l'Audit LiZbiZ
% ==============================================================================